\documentclass{article}
\usepackage{geometry}

\geometry{a4paper, margin=1in}

\title{Snapshot Objectives For EagleNav}
\author{Members of CS 3338-02 Spring '26 Group 5 @ CSULA}
\date{May 2026}

\begin{document}

\maketitle
\subsection{Disclaimer}
This document does not represent the snapshot objectives of the actual EagleNav project. Rather, this is a hypothetical document for CS 3338-02 Spring '26's final project.

\section{Introduction}
EagleNav is a mobile application designed to help users navigate the CSULA campus using augmented reality. The system provides visual directions, voice guidance, and basic navigation support to improve accessibility and ease of movement across campus.

\

To ensure that EagleNav meets the requirements outlined in the Software Requirements Specification (SRS), we have created snapshot objectives that the project is expected to meet. This document contains these snapshot objectives, with each snapshot consisting of criterion that helps validate that software both meets specifications and that implementation as described in the Software Design Document (SDD) works. These can range from core-functionalities, to niche/obscure features.

\section{Snapshot 1}
\begin{itemize}
\item Define requirements
\item Create design documents
\end{itemize}

\section{Snapshot 2}
\begin{itemize}
\item Build basic features
\item Setup UI
\end{itemize}

\section{Snapshot 3}
\begin{itemize}
\item Add AR and voice features
\item Improve system
\end{itemize}

\section{Snapshot 4}
\begin{itemize}
\item Finalize project
\item Testing and fixes
\end{itemize}

\section{Conclusion}
The snapshot objectives outline the main stages of development for EagleNav. Completing each snapshot ensures that the system is properly designed, implemented, and tested before final delivery.
\end{document}
